\subsection{The vanishing first weight and a detected higher integral weight}

\begin{statement}{A.183}{Proved}
{The evaluated norm obstruction has zero first weight and a nonzero finite integral weight}
Use the actual pair and coefficient sequence of A.181:
\[
 \begin{gathered}
 k=\Q(\alpha),\quad\alpha^3=3,\quad d=649,\quad
 A=d^8\alpha,\quad U=\alpha^2-2,\quad B=U^3,\\
 0\longrightarrow D_A\xrightarrow{T}V_A
       \xrightarrow{\epsilon}\Z_3(1)\longrightarrow0,
 \qquad \mathfrak o=\delta\kappa_3(B)\ne0.
 \end{gathered}
\tag{A.183.1}
\]
The actual quotient by the first power of the derived coordinate
has the following action and obstruction image:
\[
 D_A/TD_A=\Z_3(2),\qquad
 \mathfrak o\longmapsto
 \kappa_3(A)\cup\kappa_3(B)=0
       \quad\hbox{in }H^2_{\rm cts}(G_k,\Z_3(2)).
\tag{A.183.2}
\]
More precisely, let \(F=k(\zeta_3)\), and let \(S_F\) consist of
its primes above \(2,3,11,59\) and its complex places.  Then
\[
 3H^2_{\rm cts}(G_{F,S_F},\Z_3(2))=0,
 \qquad
 3H^2_{\rm cts}(G_{k,S_k},\Z_3(2))=0,
\tag{A.183.3}
\]
where \(S_k\) is the set below \(S_F\).

For \(s\ge1\) put \(D_{A,s}=D_A/T^sD_A\), with the quotient
arithmetic action.  There is an actual isomorphism
\[
 H^2_{\rm cts}(G_k,D_A)
       \xrightarrow{\ \sim\ }
 \varprojlim_s H^2_{\rm cts}(G_k,D_{A,s}).
\tag{A.183.4}
\]
In fact \(\mathfrak o_{14}\ne0\), so the least detecting index satisfies
\(2\le s\le14\).
Every nonzero image \(\mathfrak o_s\) belongs to the following
explicitly specified subgroup:
\[
\begin{aligned}
 \mathcal P_s&=\ker\left(H^2_{\rm cts}(G_k,D_{A,s})
       \longrightarrow\varprojlim_cH^2(G_k,D_{A,s}/3^cD_{A,s})\right),\\
 \mathcal P_s&=\operatorname{im}\left(
 \varprojlim\nolimits_c^1H^1(G_k,D_{A,s}/3^cD_{A,s})
       \longrightarrow H^2_{\rm cts}(G_k,D_{A,s})\right).
\end{aligned}
\tag{A.183.5}
\]
The map from the derived limit in A.183.5 is injective.
The exact least detecting index is proved to be two in A.195 below.
\end{statement}

\begin{proof}
\emph{The literal quotient action and the cup sign.}
Write \(X=1+T\), \(\chi=\chi_3\), and choose the continuous
Kummer cocycles \(a,b,u\) with \(b=3u\), for \(A,B,U\).
The actions in A.181.17 are
\[
\begin{aligned}
 g v&=\chi(g)X^{a(g)}\vartheta_g(v),\qquad
                    \vartheta_g(X)=X^{\chi(g)},\\
 g h&=\chi(g)X^{a(g)}
       \frac{X^{\chi(g)}-1}{X-1}\vartheta_g(h)
                    \quad(h\in D_A).
\end{aligned}
\tag{A.183.6}
\]
At \(T=0\), the multiplier in the second line is
\(\chi(g)^2\).  This gives the asserted equivariant
identification with \(\Z_3(2)\), including its chosen basis.
Use the inhomogeneous differential
\(\delta w(g,h)=g w(h)-w(gh)+w(g)\).
The constant-polynomial lift \(w(g)=b(g)\in V_A\) has
\[
 \left(\frac{\delta w(g,h)}{T}\right)\bmod T
     =a(g)\chi(g)b(h).
\tag{A.183.7}
\]
Indeed \(X^{a(g)}=1+a(g)T\pmod {T^2}\), and the constant
coefficient cancels by the cocycle identity for \(b\).
Thus the coefficient map sends \(\delta[b]\) to the cup
with the positive sign displayed in A.183.2.
It sends the same class to \(3(a\cup u)\).

All these Tate modules and Kummer cocycles factor through
\(G_{k,S_k}\).  At a finite place away from three, adjoining
prime-to-residue-characteristic roots of a unit is unramified:
its roots over the residue algebra lift by the invertible
derivative.  The valuations of \(A\) are supported above
\(3,11,59\), and \(U\) is a global unit.  This proves the
factorization on the actual root systems and hence on their
compatible cocycles.

\emph{The base ring with only three inverted.}
Set
\[
 X_0=\operatorname{Spec}\mathcal O_F[1/3],\qquad
 X=\operatorname{Spec}\mathcal O_F[1/(2\cdot3\cdot11\cdot59)],
 \qquad M_c=\Z/3^c(2).
\tag{A.183.8}
\]
The field \(F\) is exactly the field \(L_1\) in A.138.
That result proves \(\operatorname{Cl}(F)\{3\}=0\), and its actual local
holomorph proof gives precisely one prime above three.
The field is totally imaginary.  Hence
\(\operatorname{Pic}(X_0)/3=0\).
The global Brauer invariant sequence, in the form
\cite[II, Proposition 2.1]{BIR_MilneADT}, injects
\(H^2_{\acute et}(X_0,\mathbf G_m)\) into the Brauer groups
at its omitted places.  There is just one finite omitted
place and all archimedean ones are complex.  Its invariant
must sum to zero, so
\(H^2_{\acute et}(X_0,\mathbf G_m)=0\).
The Kummer sequence therefore gives
\[
 H^2_{\acute et}(X_0,\mu_3)=0.
\tag{A.183.9}
\]
Choose \(\zeta_3\in F\) to identify \(M_1\) with \(\mu_3\).
The exact coefficient sequences
\(0\to M_{c-1}\to M_c\to M_1\to0\)
and their long exact cohomology sequences now prove
\(H^2_{\acute et}(X_0,M_c)=0\) for every \(c\).

We specify the group-to-scheme comparison used here.
For a finite locally constant sheaf \(M\) on either ring in
A.183.8, the natural map
\[
 H^2(G_{F,S},M)\longrightarrow H^2_{\acute et}(X_S,M)
\tag{A.183.10}
\]
is injective.  Take the pro-cover formed by all finite
connected \(\acute etale\) covers of \(X_S\), with its actual
group \(G_{F,S}\).  After a finite cover the coefficient sheaf
is constant.  Every first-cohomology class on any further
finite stage is a finite \(\acute etale\) torsor; its pullback
to that torsor, and then to a finite Galois cover dominating
it, is zero.  Thus the direct-limit first cohomology on the
pro-cover is zero.  The five-term Hochschild--Serre sequence
gives A.183.10.  This is the degree-two portion of the
comparison proved in
\cite[II, Proposition 2.9 and its proof]{BIR_MilneADT};
the argument here uses only the stated torsor-killing map.

The first cohomology groups with each finite coefficient are
finite.  For group cohomology this is
\cite[I, Corollary 4.15]{BIR_MilneADT}; for the schemes it also
follows directly from their description by finite \(\acute etale\)
torsors and the bounded-degree finiteness theorem for number
fields unramified outside a finite set.
Both first-cohomology inverse systems therefore have zero
\(\varprojlim^1\).  The continuous inverse-limit comparison
\cite[(1.6), (2.1), Theorem (2.2), and (3.1)]{BIR_Jannsen1988}
shows that A.183.10 remains injective for \(\Z_3(2)\) and that
\[
 H^2_{\acute et,\mathrm{cts}}(X_0,\Z_3(2))=0.
\tag{A.183.11}
\]
This proves the needed integral assertion before any primes
are added to the omitted set.

\emph{The exact residue groups at the added primes.}
At each removed point \(v\mid2\cdot11\cdot59\), three is
invertible in the strict henselian discrete valuation ring.
Its generic tame pro-three inertia is \(\Z_3(1)\); wild
inertia and the other tame factors have zero positive-degree
three-primary cohomology.  The local pair sequence therefore
identifies the supported coefficient complex with
\(M_c(-1)[-2]\).  This gives, with the actual arithmetic
Frobenius and its action on the Tate basis, the localization
sequence
\[
 H^2_{\acute et,\mathrm{cts}}(X_0,\Z_3(2))
 \longrightarrow H^2_{\acute et,\mathrm{cts}}(X,\Z_3(2))
 \xrightarrow{\partial}
 \bigoplus_{v\mid2\cdot11\cdot59}
                     H^1(\mathbf F_v,\Z_3(1)).
\tag{A.183.12}
\]
Here the shift and twist come from the displayed tame-inertia
calculation.  The support and inverse-limit maps are the
ones in \cite[(3.6), (3.8), and (3.15)--(3.16)]{BIR_Jannsen1988}.
There are finitely many points, so these finite support sums
commute with the inverse limit.  Its finite first cohomology
again removes every possible \(\varprojlim^1\) term used here.
By A.183.11, \(\partial\) is injective.

The following table records every added point, grouped by
its rational prime.  The quadratic factors are irreducible
in the indicated residue fields.
\[
\begin{array}{c|c|c|c|c}
 q&X^3-3\pmod q&\#\{v\mid q\}\text{ in }F
       &|\mathbf F_v|&v_3(|\mathbf F_v|-1)\\ \hline
 2&(X-1)(X^2+X+1)&3&4&1\\
 11&(X-9)(X^2-2X+4)&3&121&1\\
 59&(X-49)(X^2-10X+41)&3&3481&1
\end{array}
\tag{A.183.13}
\]
Indeed each factorization is squarefree with cycle type
\((1)(2)\); in the actual \(S_3\) splitting field its
Frobenius has order two, hence there are three primes of
residue degree two.  For a finite field of size \(Q\), the
continuous procyclic cochain calculation is
\[
 H^1(\mathbf F_Q,\Z_3(1))
       =\operatorname{coker}\bigl(Q-1:\Z_3\to\Z_3\bigr)
       =\Z_3/(Q-1)\Z_3.
\tag{A.183.14}
\]
One may obtain this integral map by taking the inverse limit
of the two-term finite-coefficient complexes; their terms
are onto and their finite zeroth cohomology has zero derived
limit.  Each of the nine summands in A.183.12 is consequently
killed by three.  The injective maps A.183.10--12 prove the
first equality in A.183.3.

\emph{Restriction, corestriction, and the evaluated cup.}
The extension \(F/k\) has degree two, and the omitted prime
sets in A.183.8 are inverse images of those over \(k\).
Thus the actual restriction and corestriction maps satisfy
\[
 \operatorname{cor}_{F/k}\operatorname{res}_{F/k}=2
       \quad\hbox{on }H^2_{\rm cts}(G_{k,S_k},\Z_3(2)).
\tag{A.183.15}
\]
This holds on finite cochains and commutes with all coefficient
maps, hence also holds integrally.  Multiplication by two is
an automorphism of the coefficient module and its cohomology.
The exponent-three assertion over \(F\) therefore proves the
one over \(k\).  In particular \(3(a\cup u)=0\) over
\(G_{k,S_k}\), and its inflation to \(G_k\) is zero.
Together with A.183.7 this proves A.183.2.

\emph{The integral weight limit over this actual base field.}
The filtration of \(D_{A,s}\) by powers of \(T\) has successive
quotients
\(\Z_3(2),\Z_3(3),\ldots,\Z_3(s+1)\), directly from A.183.6.
It is unramified outside \(S_k\).  At a place \(v\notin S_k\)
with residue cardinality \(Q\), restriction of a continuous
one-cocycle to inertia lies in
\[
 H^1(I_v,D_{A,s})^{\operatorname{Frob}_v}
       =D_{A,s}(-1)^{\operatorname{Frob}_v}=0.
\tag{A.183.16}
\]
For the first equality, inertia acts trivially and its pro-three
tame quotient is \(\Z_3(1)\).  For the second, the matrix
\(\operatorname{Frob}_v-1\) on this finite free lattice is
triangular with nonzero diagonal entries
\(Q-1,Q^2-1,\ldots,Q^s-1\); its determinant is nonzero and
therefore its kernel is zero.  A cocycle consequently vanishes
on these inertia groups and their closed normal closure.
Inflation gives the literal equality
\[
 H^1_{\rm cts}(G_k,D_{A,s})
       =H^1_{\rm cts}(G_{k,S_k},D_{A,s}).
\tag{A.183.17}
\]
The right hand side is an inverse limit of finite first
cohomology groups: its finite zeroth-cohomology system is
Mittag--Leffler.  It is compact, and its quotient by three
injects into the finite group
\(H^1(G_{k,S_k},D_{A,s}/3D_{A,s})\).
Successive approximation modulo powers of three now gives
finite generation over \(\Z_3\).  In particular all transition
maps in \(s\) are continuous maps of compact groups.

The derived inverse limit of a countable system of compact
abelian groups is zero: every finite initial system of equations
\(z_s-\pi_s z_{s+1}=c_s\) can be solved backwards, and their
closed solution sets have the finite intersection property in
the compact product.  Also \(D_A=\varprojlim_sD_{A,s}\) with
its actual compact topology, and the quotient maps have
continuous linear sections.  Continuous cochains are therefore
degreewise onto.  Their Milnor exact sequence and A.183.17
prove A.183.4.  This repeats the actual A.81.8 maps over
\(k\), with the finite-support argument supplied here.

Since A.181 proves \(\mathfrak o\ne0\), some finite integral
image is nonzero; A.183.2 puts its index above one.
For each fixed \(s\), apply the same cochain inverse-limit
sequence to the coefficients \(D_{A,s}/3^cD_{A,s}\).
Their set sections give the exact injection and kernel in
A.183.5.  A.181.23 supplies the actual equivariant maps from
its root-torsor quotients to every such finite coefficient
quotient, and sends the obstruction to zero there.
Thus every detecting \(\mathfrak o_s\) lies in the stated
nonzero coefficient-power derived-limit subgroup.
The compactness in A.183.17 concerns the integral modules.
No finiteness claim for \(H^1(G_k,D_{A,s}/3^cD_{A,s})\) is
used: these finite-coefficient classes can ramify at primes
outside the fixed set, and their displayed derived limit is
precisely where the nonzero image has been proved to lie.

\emph{An explicit detecting integral quotient.}
All coefficient sequences below first use \(G_{k,S_k}\).
Put
\[
 L_2=\frac{(1+T)^9-1}{T},\qquad
 J_2=D_A/L_2D_A,\qquad
 V_2=\Z_3[\operatorname{Roots}(Z^9-A)](1).
\tag{A.183.18}
\]
These are the literal level-two root-torsor coefficients in
A.181.18--19.  Shapiro identifies the augmentation map on
first cohomology with the completed \(S\)-unit norm from
\(K_3\) to \(k\).  This follows also from the finite Kummer
sequences: their class-group terms have zero inverse Tate
module because the class groups are finite.  Hence the
connecting map embeds
\[
 C=\widehat{\mathcal O_{k,S_k}^{\times}}_3/
       N_{K_3/k}\widehat{\mathcal O_{K_3,S_3}^{\times}}_3
       \hookrightarrow
 H^2_{\rm cts}(G_{k,S_k},J_2).
\tag{A.183.19}
\]
Scalar norms kill \(C\) by nine.  The image of \(U\) has
exact order nine: its ninth power is the scalar norm of
\(U\), whereas its cube is not a completed \(S\)-unit norm
by A.181.3.  Write \(\delta_S\) for this connecting map.
Then \(\delta_S(B)=3\delta_S(U)\ne0\).

The group \(H^2_{\rm cts}(G_{k,S_k},V_2)\) is torsion-free.
Indeed Shapiro sends it to
\(H^2_{\rm cts}(G_{K_3,S_3},\Z_3(1))\), and A.183.10,
with the identical proof over \(K_3\), injects this into
integral \(\acute etale\) cohomology on
\(\operatorname{Spec}\mathcal O_{K_3,S_3}\).
The three-primary \(S\)-class group is zero by A.177 and
A.179.  The finite Kummer sequences identify its degree-two
cohomology with \(\operatorname{Br}(\mathcal O_{K_3,S_3})[3^c]\).
The invariant sequence identifies the inverse limit, with
its multiplication-by-three transitions, with the kernel
of the sum map \(\Z_3^{S_{3,\mathrm{fin}}}\to\Z_3\).
Real Brauer groups have no three-primary part.
This kernel is a finite free \(\Z_3\)-module, proving the
claimed torsion-freeness through the actual injection.

Now \(\delta_S(B)\) is not divisible by nine in
\(H^2(G_{k,S_k},J_2)\).  If \(9x=\delta_S(B)\), the image
of \(x\) in the torsion-free group just identified is killed
by nine, hence zero.  Exactness puts \(x\) in \(C\), where
nine kills every element, contradicting \(\delta_S(B)\ne0\).
The coefficient sequence for multiplication by nine thus
sends \(\delta_S(B)\) to a nonzero class in
\(H^2(G_{k,S_k},J_2/9J_2)\).

The polynomial quotient needed to reach this class is exact:
\[
 \begin{gathered}
 L_2\equiv T^8+3T^5+3T^2\pmod9,\qquad
 T^{14}\in(9,L_2),\\
 D_A/T^{14}D_A\longrightarrow J_2/9J_2.
 \end{gathered}
\tag{A.183.20}
\]
For example the relation gives
\(T^{11}\equiv-3T^5\) and then \(T^{14}=0\) modulo
\((9,L_2)\).  These are equivariant coefficient quotients
of the same module \(D_A\).  The finite-set obstruction
\(\mathfrak o_{S,14}\) therefore maps to the nonzero class
just constructed.

Finally the inflation
\[
 H^2_{\rm cts}(G_{k,S_k},D_{A,14})
       \longrightarrow H^2_{\rm cts}(G_k,D_{A,14})
\tag{A.183.21}
\]
is injective.  Here is a direct cochain proof, which also
specifies the continuous map.  Let \(N\) be the kernel of
\(G_k\to G_{k,S_k}\), and let an inflated normalized
2-cocycle be the coboundary of a continuous 1-cochain \(b\).
Its restriction to \(N\) is a homomorphism to \(D_{A,14}\).
The cocycle identities with one entry in \(N\) show that
this homomorphism is \(G_k\)-equivariant.  Its restriction
to every outside-\(S_k\) inertia group is zero by A.183.16.
Those inertia groups normally generate \(N\), so it is zero
on \(N\).  The same identities give \(b(gn)=b(g)\).
A continuous map constant on the fibres of the compact
quotient descends continuously, so \(b\) is already a
cochain on \(G_{k,S_k}\).  This proves A.183.21 and hence
\(\mathfrak o_{14}\ne0\).
The nonzero finite-coefficient test used in this paragraph
is over \(G_{k,S_k}\); its inflation to \(G_k\) is zero
by A.181.23.  The maps A.183.20--21 exhibit exactly how
that fact coexists with the nonzero integral image.
\end{proof}

\paragraph{Exact finite certificate.}
\path{code/two_inert_first_weight_certificate.py} and its retained
output and data check every factorization and residue size in
A.183.13, the finite procyclic coefficient and transfer maps,
the constant and linear terms of the actual action in
A.183.6--7, and the full integer division of \(T^{14}\) by
\(L_2\) with remainder divisible by nine.  The arithmetic cohomology and infinite-limit
assertions are established by the maps and arguments above.
